\chapter{一致收敛的保持定理}
\label{chap:uniform-convergence-preservation-solutions}

\begin{solution}{uniform-limit-continuity-epsilon}
给定 \(\varepsilon>0\)，先由一致收敛取 \(N\)，使
\[
 |f_N(x)-f(x)|<\varepsilon/3\qquad(x\in E).
\]
再由 \(f_N\) 在 \(x_0\) 连续取 \(\delta>0\)，使
\(x\in E\)、\(|x-x_0|<\delta\) 时
\[
 |f_N(x)-f_N(x_0)|<\varepsilon/3.
\]
将 \(f_N(x)\) 与 \(f_N(x_0)\) 插入 \(f(x)-f(x_0)\)，三角不等式给出
\(|f(x)-f(x_0)|<\varepsilon\)。这就是 \(f\) 在 \(x_0\) 的连续性。
\end{solution}

\begin{solution}{local-uniform-geometric-series}
记 \(S_n(x)=\sum_{k=0}^nx^k\)。在 \([-r,r]\)、\(r<1\) 上，
\[
 \left|S_n(x)-\frac1{1-x}\right|
 =\frac{|x|^{n+1}}{|1-x|}
 \le\frac{r^{n+1}}{1-r}\to0.
\]
任意紧子区间 \(K\subset(-1,1)\) 都包含在某个这样的 \([-r,r]\) 中，
故收敛局部一致。对每个固定 \(n\)，当 \(x\to1-\) 时上式的误差趋于
\(+\infty\)，所以误差在 \((-1,1)\) 上甚至无界，不可能全局一致收敛。
\end{solution}

\begin{solution}{uniform-integrability-darboux}
设分割 \(P:a=x_0<\cdots<x_m=b\)，\(I_j=[x_{j-1},x_j]\)。由
\[
 f_N(x)-\eta\le f(x)\le f_N(x)+\eta,\qquad
 \eta=\norm{f-f_N}_\infty,
\]
得到
\[
 \sup_{I_j}f\le\sup_{I_j}f_N+\eta,\qquad
 \inf_{I_j}f\ge\inf_{I_j}f_N-\eta.
\]
所以
\[
 \omega(f;I_j)\le\omega(f_N;I_j)+2\eta.
\]
乘以 \(\Delta x_j\) 后求和，
\[
 U(f,P)-L(f,P)
 \le U(f_N,P)-L(f_N,P)+2\eta(b-a).
\]
一致收敛使一个固定的充分靠后的 \(f_N\) 在所有子区间上同时近似 \(f\)，
所以无需为不同子区间选择不同指标。
\end{solution}

\begin{solution}{uniform-integral-varying-upper-limit}
对任意 \(a\le u\le v\le b\)，
\[
 \left|\int_u^v(f_n-f)\right|
 \le(v-u)\norm{f_n-f}_\infty
 \le(b-a)\norm{f_n-f}_\infty.
\]
右端与 \(u,v\) 无关，故对所有区间取上确界仍得到同一上界。
\end{solution}

\begin{solution}{pointwise-integral-spike}
在 \(x=1/n\) 处，左侧公式给出零，与右侧零值相接，所以 \(s_n\) 连续。
固定 \(x>0\) 时最终 \(x>1/n\)，故函数值为零；在 \(x=0\) 处也恒为零。
令 \(u=nx\)，则在支集上
\[
 s_n(x)=n\,u(1-u),
\]
其最大值 \(n/4\) 在 \(u=1/2\) 取得，并且
\[
 \int_0^1s_n(x)\,dx
 =\int_0^1u(1-u)\,du=\frac16.
\]
若 \(s_n\) 一致趋于零，则积分保持定理会迫使积分趋于零，与 \(1/6\)
矛盾；也可直接由 \(\norm{s_n}_\infty=n/4\) 看出。
\end{solution}

\begin{solution}{uniform-integral-with-moving-interval}
把差写成
\[
 \int_{a_n}^{b_n}(f_n-f)
 +\left(\int_{a_n}^{b_n}f-\int_\alpha^\beta f\right).
\]
第一项绝对值不超过
\((b-a)\norm{f_n-f}_\infty\)。函数 \(f\) 在闭区间上有界，第二项可由
有向积分的可加性估计为
\[
 \left|\int_{a_n}^{\alpha}f\right|
 +\left|\int_{\beta}^{b_n}f\right|
 \le\norm f_\infty\bigl(|a_n-\alpha|+|b_n-\beta|\bigr).
\]
两部分都趋于零。
\end{solution}

\begin{solution}{derivative-theorem-anchor}
导数只决定函数到一个加法常数。一个点处的收敛控制这个常数，而中值定理
把导数的共同控制传到其他点。若 \(f_n(x)=x+n\)，则
\(f_n'\equiv1\) 一致收敛，但对每个 \(x\)，\(f_n(x)\to+\infty\)，
函数列没有有限极限。因此完全删除锚点条件是不行的。
\end{solution}

\begin{solution}{derivative-theorem-uniform-function-convergence}
对每个 \(x\in[a,b]\)，中值定理给出
\[
 |(f_m-f_n)(x)-(f_m-f_n)(x_0)|
 \le |x-x_0|\norm{f_m'-f_n'}_\infty.
\]
再用 \(|x-x_0|\le b-a\)，得到
\[
 |f_m(x)-f_n(x)|
 \le |f_m(x_0)-f_n(x_0)|
 +(b-a)\norm{f_m'-f_n'}_\infty.
\]
对 \(x\) 取上确界即为所求。
\end{solution}

\begin{solution}{uniform-functions-pointwise-derivatives}
令 \(t=n|x|\)。则
\[
 |f_n(x)|=\frac1n\frac{t}{1+t^2}\le\frac1{2n},
\]
等号在 \(|x|=1/n\) 处取得，故 \(f_n\rightrightarrows0\)。求导得
\[
 f_n'(x)=\frac{1-n^2x^2}{(1+n^2x^2)^2}.
\]
若 \(x=0\)，该值恒为 \(1\)；若 \(x\ne0\)，分子是 \(O(n^2)\)，
分母是 \(O(n^4)\)，故极限为零。因此导数逐点极限在原点取 \(1\)、其余
点取零，不等于零函数的导数。
\end{solution}

\begin{solution}{uniform-derivative-no-pointwise-function-given}
由习题 \ref{ex:derivative-theorem-uniform-function-convergence} 的估计，
右端在 \(m,n\to\infty\) 时趋于零：第一项因锚点值收敛而趋零，第二项
因导数列一致收敛而趋零。因此 \((f_n)\) 是一致 Cauchy 列。实值函数列
的一致 Cauchy 准则直接给出一致收敛，故不需另设逐点收敛。
\end{solution}

\begin{solution}{termwise-derivative-series-form}
令
\[
 S_N=\sum_{n=1}^Nu_n.
\]
则 \(S_N(x_0)\) 收敛，且
\[
 S_N'=\sum_{n=1}^Nu_n'
\]
一致收敛。对函数列 \((S_N)\) 使用逐项求导定理，得到它一致收敛于某个
\(S\)，并且
\[
 S'=\lim_{N\to\infty}S_N'=\sum_{n=1}^\infty u_n'.
\]
这正是函数项级数的逐项求导结论。
\end{solution}

\begin{solution}{derivative-limit-darboux}
正文定理已经证明 \(g=f'\)。第一卷的 Darboux 定理说明任何导函数都具有
介值性质，所以 \(g\) 具有 Darboux 性。这个结论来自 \(g\) 是某个函数
的导数，而不是来自一般的一致极限性质。
\end{solution}

\begin{solution}{dini-closed-set-proof}
设 \(f_n\downarrow f\)，并令
\[
 F_n=\{x\in K:f_n(x)-f(x)\ge\varepsilon\}.
\]
每个 \(F_n\) 是 \(K\) 中的闭集，且 \(F_{n+1}\subseteq F_n\)。
逐点收敛说明 \(\bigcap_nF_n=\varnothing\)。若所有 \(F_n\) 都非空，
紧集中的嵌套非空闭集定理会推出交集非空，矛盾。因此某个 \(F_N\) 为空。
对 \(n\ge N\)，\(F_n\subseteq F_N\)，故
\(0\le f_n-f<\varepsilon\) 在 \(K\) 上处处成立。
\end{solution}

\begin{solution}{dini-compactness-counterexample}
对固定 \(x\ge0\)，分母随 \(n\) 增大，所以
\[
 q_n(x)=\frac{x}{n+x}\downarrow0.
\]
各函数和极限都连续。然而 \(q_n(x)\to1\) 当 \(x\to\infty\)，故
\[
 \sup_{x\ge0}q_n(x)=1
\]
对每个 \(n\) 成立，收敛不一致。定义域不是紧集。
\end{solution}

\begin{solution}{dini-continuous-limit-counterexample}
在 \([0,1]\) 上，\(x^n\) 对 \(n\) 逐点单调递减，且每项连续。逐点极限
在 \([0,1)\) 上为零、在 \(1\) 处为一，因而不连续。又
\(\sup_{[0,1]}|x^n-f(x)|=1\)（上确界由趋近 \(1\) 的左侧点给出），
所以收敛不一致。
\end{solution}

\begin{solution}{dini-monotonicity-counterexample}
可取上一章的移动峰
\[
 p_n(x)=\max\{1-3n|x-1/n|,0\}.
\]
它连续且逐点趋于零，但 \(\norm{p_n}_\infty=1\)。对某些固定
\(x>0\)，当移动支集经过 \(x\) 时函数值先从零升高再回到零，因此序列
\((p_n(x))\) 一般不关于 \(n\) 单调。更直接地，也可把峰中心安排为在区间
内依次移动，使同一点被若干峰先后经过。
\end{solution}

\begin{solution}{dini-uniform-supremum-proof}
令 \(g_n=f_n-f\ge0\)，并取 \(x_n\in K\) 使
\[
 M_n=\max_Kg_n=g_n(x_n).
\]
因 \(g_{n+1}\le g_n\)，数列 \(M_n\) 单调递减，设其极限为 \(M\ge0\)。
由紧致性，\((x_n)\) 有子列 \(x_{n_k}\to x_*\in K\)。固定 \(m\)，
充分大的 \(k\) 有 \(n_k\ge m\)，从而
\[
 M_{n_k}=g_{n_k}(x_{n_k})\le g_m(x_{n_k})\to g_m(x_*).
\]
故 \(M\le g_m(x_*)\) 对每个 \(m\) 成立。令 \(m\to\infty\)，逐点收敛
给出 \(M=0\)。紧致性正用于从最大点列抽取仍落在 \(K\) 中的收敛子列。
\end{solution}

\begin{solution}{uniform-limit-homeomorphism-inverses}
给定 \(\varepsilon>0\)。因 \(f^{-1}\) 在紧区间 \([c,d]\) 上一致连续，
存在 \(\eta>0\)，使 \(|u-v|<\eta\) 时
\[
 |f^{-1}(u)-f^{-1}(v)|<\varepsilon.
\]
取 \(N\) 使 \(\norm{f_n-f}_\infty<\eta\)。对任意 \(y\in[c,d]\)，令
\(x=f_n^{-1}(y)\)。则
\[
 |f(x)-y|=|f(x)-f_n(x)|<\eta,
\]
于是
\[
 |f_n^{-1}(y)-f^{-1}(y)|
 =|f^{-1}(f(x))-f^{-1}(y)|<\varepsilon.
\]
估计对所有 \(y\) 同时成立，故逆函数列一致收敛。
\end{solution}
